\documentclass[../DoAn.tex]{subfiles}
\begin{document}

\begin{center}
    \Large{\textbf{ABSTRACT}}\\
\end{center}
\vspace{1cm}
There is currently a great demand for in-depth legal advice on marriage and family matters due to the potential legal complications that can arise during various stages of a marital relationship. When seeking advice, users often require answers based on precise legal arguments, specifically cited with clear clause numbers and warnings about validity and conflicts (if any). Furthermore, due to the hallucinations of artificial intelligence, users also need to observe and verify the chatbot's reasoning and question processing before receiving a final answer to fully trust the chatbot's response.

In the Vietnamese market, there are already tools such as AITracuuluat, AI Law, the National Legal Portal's chatbot, and popular large-scale language models like Gemini. However, practical surveys show that these products still have limitations, such as sometimes inaccurate legal basis, a lack of mechanisms to warn about the validity of legal documents and legal conflicts when a situation is governed by multiple overlapping regulations, and finally, a lack of a transparent mechanism to display the chatbot's processing and reasoning flow before providing a final answer.

To overcome these limitations, the overall solution of this project is to build a chatbot application capable of automatically advising on Marriage and Family Law and transparently displaying the reasoning flow before providing a final answer. The project chooses to build the application using the Agentic GraphRAG architecture, combining Knowledge Graph with Agentic AI. Knowledge graphs allow for the modeling of legal entities and relationships between documents (references, instructions, amendments, replacements, repeals, conflicts), supporting the retrieval of accurate and valid legal grounds and providing warnings about validity and conflicts; Agentic AI optimizes querying on knowledge graphs by separating retrievers corresponding to question types and using agents to decide which retriever(s) to use to retrieve data for generating answers.

% The main technologies used for building the chatbot application include the Chainlit library to build a chatbot interface that visualizes the processing flow; vis-network to visually plot the retrieved knowledge graph; GPT-40 API for Router Agent and Retriever Agent, GPT-4.1 for answer generation; Neo4j for storing and retrieving knowledge graphs; and PostgreSQL for storing user information, conversation sessions, and agent processing flows. Chainlit for the chat interface. The knowledge graph was built from the 2014 Marriage and Family Law and more than 20 related documents (Civil Code, Law on Civil Status, Law on Prevention of Domestic Violence, Decrees, Resolutions and Circulars) with detailed breakdown down to the point level of the law.
The chatbot was evaluated using a combination of manual black-box testing and automated evaluation on a test case of 200 question-answer pairs collected from the Legal Q\&A section on the Legal Library, using the RAGAS library and specific indicators of legal basis. The results showed that the system achieved Faithfulness 88.2\%, Context Recall 83.5\% on RAGAS, and Legal Citation Accuracy 83.7\%, ensuring accuracy and reliability in answers for users

\end{document}